\chapter{Fragment Orbitals and Polyatomic Molecules}\label{ch:b2:fragment-orbitals}

The Lewis structure of water shows two equivalent O--H bonds and two
equivalent lone pairs. Yet when water is ionised by ultraviolet light, the
electrons come out with four different energies, not two; methane, with four
equivalent C--H bonds, shows two. Electrons in a molecule do not sit in bonds
between pairs of atoms: they occupy orbitals spread over the whole molecule.
This chapter builds those orbitals from the orbitals of \omterm{def:b2:fragment-orbitals:fragment}{fragments}, explains
why water is bent and ethene is flat, and why a carbocation prefers to have
alkyl groups around it.

\begin{recall}
\Cref{ch:b2:diatomic-mos}: \omterm{def:b2:diatomic-mos:lcao}{LCAO}, overlap and symmetry, two interacting levels,
bonding and \omterm{def:b2:diatomic-mos:bonding}{antibonding orbitals}. The Year 1 volume: VSEPR shapes, hybridisation
as a description, resonance, carbocations and their order of stability (the
interaction of the C--H bonds with the empty orbital was announced there).
\end{recall}

\section{The fragment method}

\begin{definition}[Fragments]\label{def:b2:fragment-orbitals:fragment}
A \emph{fragment}\index{fragment} is a group of atoms of a molecule considered on
its own (an atom, a pair of hydrogens, a \ce{CH2} group), and a \emph{fragment
orbital}\index{fragment orbital} one of its orbitals. The orbitals of the
molecule are built by letting the orbitals of two fragments interact, two by
two, as the atomic orbitals of a diatomic molecule.
\end{definition}

\begin{definition}[Symmetry-adapted combination]\label{def:b2:fragment-orbitals:symmetry-adapted}
A \emph{symmetry-adapted combination}\index{symmetry-adapted combination} of
equivalent atomic orbitals is a combination that each symmetry operation of the
molecule (a reflection in a mirror plane, a rotation about an axis) either
leaves unchanged or turns into its opposite, or, for degenerate sets, transforms
into combinations of the same set.
\end{definition}

\begin{proposition}[Two hydrogens]\label{prop:b2:fragment-orbitals:h2-pair}
For two equivalent hydrogens with 1s orbitals $h_1$ and $h_2$, the
\omterm{def:b2:fragment-orbitals:symmetry-adapted}{symmetry-adapted combinations} are $h_1 + h_2$ and $h_1 - h_2$: the first is
unchanged, the second changes sign, under any operation that exchanges the two
atoms.
\end{proposition}

\begin{proof}
An operation that exchanges the atoms turns $h_1$ into $h_2$ and $h_2$ into
$h_1$: $h_1 + h_2$ is unchanged and $h_1 - h_2$ becomes $h_2 - h_1 = -(h_1 - h_2)$.
\end{proof}

\begin{proposition}[Three and four hydrogens]\label{prop:b2:fragment-orbitals:h4}
For four hydrogens at the corners of a tetrahedron, the combinations are one
totally symmetric combination $h_1 + h_2 + h_3 + h_4$ and a set of three degenerate
combinations with one nodal plane each, shaped like p orbitals. For three
hydrogens of a triangle (as in ammonia), one symmetric combination and a
degenerate pair.
\end{proposition}

\admitted

The shapes of the degenerate combinations, and their names (a, e, t), come from
group theory, developed in the Year 3 volume; here the names are only labels.

\begin{proposition}[Only like interacts with like]\label{prop:b2:fragment-orbitals:interaction-rules}
A \omterm{def:b2:fragment-orbitals:fragment}{fragment orbital} interacts only with orbitals of the other \omterm{def:b2:fragment-orbitals:fragment}{fragment} that
behave in the same way under every symmetry operation of the molecule; the
strength of the interaction grows with the overlap and falls with the energy
gap.
\end{proposition}

\begin{proof}
If two orbitals behave differently under some reflection, one is symmetric and
the other antisymmetric with respect to that plane, and their overlap and
\omterm{def:b2:diatomic-mos:integrals}{resonance integral} vanish (\cref{prop:b2:diatomic-mos:symmetry}). The dependence on
overlap and gap is that of \cref{thm:b2:diatomic-mos:heteronuclear}.
\end{proof}

\begin{method}[A fragment diagram]\label{met:b2:fragment-orbitals:fragment-diagram}
\begin{enumerate}
  \item Split the molecule into two \omterm{def:b2:fragment-orbitals:fragment}{fragments}, usually a central atom and the set
        of atoms around it.
  \item Build the \omterm{def:b2:fragment-orbitals:symmetry-adapted}{symmetry-adapted combinations} of the outer orbitals.
  \item Sort the orbitals of both \omterm{def:b2:fragment-orbitals:fragment}{fragments} by their behaviour under the mirror
        planes and axes of the molecule.
  \item Let each orbital interact with those of the same behaviour, more strongly
        when close in energy; leftovers stay nonbonding.
  \item Fill with the valence electrons and read: highest occupied level,
        bonding character, shape preferences.
\end{enumerate}
\end{method}

\section{Water}

Place water in the $yz$ plane, the $z$ axis bisecting the H--O--H angle. Two
mirror planes contain the $z$ axis: the molecular plane $yz$ and the plane $xz$
perpendicular to it. With respect to the $xz$ plane, which exchanges the two
hydrogens, $h_1 + h_2$ is symmetric and $h_1 - h_2$ antisymmetric. On oxygen, 2s and
$2\mathrm p_z$ are symmetric with respect to both planes; $2\mathrm p_y$ (in the
molecular plane, perpendicular to $z$) is antisymmetric with respect to $xz$;
$2\mathrm p_x$ is antisymmetric with respect to the molecular plane, which no
hydrogen combination is.

So $h_1 + h_2$ mixes with 2s and $2\mathrm p_z$, $h_1 - h_2$ with $2\mathrm p_y$, and
$2\mathrm p_x$ stays nonbonding. Six orbitals result, four of them filled by the eight
valence electrons.

\begin{omfigure}
\begin{tikzpicture}[font=\small]
  % O fragment (left)
  \pic (os) at (0,-2.4) {omlevel={ud}{0.7}}; \node[left] at (-0.45,-2.4) {2s};
  \pic (op1) at (-0.45,0) {omlevel={ud}{0.4}}; \pic (op2) at (0,0) {omlevel={u}{0.4}}; \pic (op3) at (0.45,0) {omlevel={u}{0.4}};
  \node[left] at (-0.7,0) {2p};
  \node at (0,-3.2) {O};
  % H2 fragment (right)
  \pic (hp) at (6,0.9) {omlevel={u}{0.7}}; \node[right] at (6.4,0.9) {$h_1 + h_2$};
  \pic (hm) at (6,1.8) {omlevel={u}{0.7}}; \node[right] at (6.4,1.8) {$h_1 - h_2$};
  \node[anchor=north west] at (6.4,0.6) {\ce{H2} fragment};
  % MOs
  \pic (m1) at (3,-2.9) {omlevel={ud}{0.7}};  \node[omsymlabel, below=0.17cm, inner sep=0.5pt] at (m1-c) {$2a_1$};
  \pic (m2) at (3,-1.55) {omlevel={ud}{0.7}};  \node[omsymlabel, below=0.17cm, inner sep=0.5pt] at (m2-c) {$1b_2$};
  \pic (m3) at (3,-0.78) {omlevel={ud}{0.7}};  \node[omsymlabel, below=0.17cm, inner sep=0.5pt] at (m3-c) {$3a_1$};
  \pic (m4) at (3,0) {omlevel={ud}{0.7}};  \node[omsymlabel, below=0.17cm, inner sep=0.5pt] at (m4-c) {$1b_1$};
  \pic (m5) at (3,2.4) {omlevel={empty}{0.7}}; \node[omsymlabel, below=0.17cm, inner sep=0.5pt] at (m5-c) {$4a_1$};
  \pic (m6) at (3,3.1) {omlevel={empty}{0.7}}; \node[omsymlabel, below=0.17cm, inner sep=0.5pt] at (m6-c) {$2b_2$};
  \draw[omcorr, densely dashed] (os-r) -- (m1-l) (op3-r) -- (m2-l) (op3-r) -- (m3-l) (op3-r) -- (m4-l) (op3-r) -- (m5-l) (op3-r) -- (m6-l);
  \draw[omcorr, densely dashed] (hp-l) -- (m1-r) (hm-l) -- (m2-r) (hp-l) -- (m3-r) (hp-l) -- (m5-r) (hm-l) -- (m6-r);
\end{tikzpicture}

{\small \omterm{def:b2:fragment-orbitals:fragment}{Fragment} diagram of water, schematic. The symmetric combination of the
hydrogens mixes with the 2s and $2\mathrm p_z$ orbitals of oxygen ($a_1$ orbitals),
the antisymmetric one with $2\mathrm p_y$ ($b_2$); $2\mathrm p_x$, perpendicular to the
molecular plane, stays nonbonding ($1b_1$, at the level of the oxygen 2p). Four filled valence
levels, four photoelectron bands; the highest, $1b_1$, is ionised at
\qty{12.62}{eV}. The labels
$a_1$, $b_1$, $b_2$ are names (from group theory, in the Year 3 volume).}
% ledger: ie:H2O
\end{omfigure}

\begin{proposition}[Why water is bent]\label{prop:b2:fragment-orbitals:bent-water}
In a linear H--O--H, two oxygen p orbitals perpendicular to the axis are
nonbonding. Bending lets one of them, the one in the molecular plane, overlap
with $h_1 + h_2$ and mix with the 2s orbital: that filled level drops, by more than
the other filled levels rise. With eight valence electrons the bent molecule is
the more stable.
\end{proposition}

\begin{proof}[Argument]
In the linear molecule ($z$ along H--O--H), $h_1 + h_2$ has zero overlap with
$2\mathrm p_x$ and $2\mathrm p_y$ (each is antisymmetric with respect to a plane
containing the axis, the combination symmetric); they are two degenerate
nonbonding pairs. On bending in the $yz$ plane, $2\mathrm p_y$ points partly towards
the hydrogens: its overlap with $h_1 + h_2$ grows from zero, and the filled level
it carries is stabilised (the $3a_1$ level of the bent molecule). The bonding
$1b_2$ level loses a little overlap and rises slightly. The net is a lowering,
largest when the descending level is filled, as it is with eight electrons. A
molecule with only four valence electrons, such as \ce{BeH2}, has that level
empty and is linear.
\end{proof}

\section{Methane, ammonia and photoelectron spectra}

In methane the four hydrogen orbitals give the symmetric combination, which
matches the 2s orbital of carbon, and a degenerate set of three, which matches
its three 2p orbitals. Four \omterm{def:b2:diatomic-mos:bonding}{bonding orbitals} are filled: one low ($a_1$) and three
degenerate at a higher energy ($t_2$). Ammonia, with three hydrogens, keeps a
filled orbital pointing away from the hydrogens on nitrogen, its lone pair, as
the highest occupied level.

\begin{omfigure}
\begin{tikzpicture}[font=\small]
  \pic (cs) at (0,-1.6) {omlevel={ud}{0.7}}; \node[left] at (-0.45,-1.6) {2s};
  \pic (cp1) at (-0.45,0) {omlevel={u}{0.4}}; \pic (cp2) at (0,0) {omlevel={u}{0.4}}; \pic (cp3) at (0.45,0) {omlevel={empty}{0.4}};
  \node[left] at (-0.7,0) {2p};
  \node at (0,-2.5) {C};
  \pic (ha) at (6,-0.8) {omlevel={ud}{0.7}}; \node[right] at (6.4,-0.8) {$a_1$};
  \pic (ht1) at (5.55,0.4) {omlevel={u}{0.4}}; \pic (ht2) at (6,0.4) {omlevel={u}{0.4}}; \pic (ht3) at (6.45,0.4) {omlevel={empty}{0.4}};
  \node[right] at (6.7,0.4) {$t_2$};
  \node at (6,-1.6) {\ce{H4} fragment};
  \pic (m1) at (3,-2.4) {omlevel={ud}{0.7}}; \node[omsymlabel, below=0.17cm, inner sep=0.5pt] at (m1-c) {$1a_1$};
  \pic (m2a) at (2.55,-0.9) {omlevel={ud}{0.4}}; \pic (m2b) at (3,-0.9) {omlevel={ud}{0.4}}; \pic (m2c) at (3.45,-0.9) {omlevel={ud}{0.4}};
  \node[omsymlabel, below=0.17cm, inner sep=0.5pt] at (m2b-c) {$1t_2$};
  \pic (m3) at (3,1.8) {omlevel={empty}{0.7}}; \node[omsymlabel, below=0.17cm, inner sep=0.5pt] at (m3-c) {$2a_1$};
  \pic (m4a) at (2.55,2.5) {omlevel={empty}{0.4}}; \pic (m4b) at (3,2.5) {omlevel={empty}{0.4}}; \pic (m4c) at (3.45,2.5) {omlevel={empty}{0.4}};
  \node[omsymlabel, below=0.17cm, inner sep=0.5pt] at (m4b-c) {$2t_2$};
  \draw[omcorr, densely dashed] (cs-r) -- (m1-l) (cp3-r) -- (m2a-l) (cs-r) -- (m3-l) (cp3-r) -- (m4a-l);
  \draw[omcorr, densely dashed] (ha-l) -- (m1-r) (ht1-l) -- (m2c-r) (ha-l) -- (m3-r) (ht1-l) -- (m4c-r);
\end{tikzpicture}

{\small \omterm{def:b2:fragment-orbitals:fragment}{Fragment} diagram of methane, schematic. The symmetric combination of the
four hydrogens matches the 2s orbital of carbon, the three degenerate ones its 2p
orbitals: two filled levels, hence two photoelectron bands, the upper one
(ionised first, at \qty{12.61}{eV}) holding six electrons.}
% ledger: ie:CH4
\end{omfigure}

\begin{definition}[Photoelectron spectroscopy]\label{def:b2:fragment-orbitals:photoelectron}
\emph{Photoelectron spectroscopy}\index{photoelectron spectroscopy} measures the
kinetic energies of the electrons ejected from molecules by monochromatic
radiation of known energy; each band of the spectrum corresponds to the removal
of an electron from one occupied level, and its position gives the ionisation
energy of that level.
\end{definition}

\begin{proposition}[Koopmans' approximation]\label{prop:b2:fragment-orbitals:koopmans}
The energy needed to remove an electron from an occupied orbital is
approximately minus the energy of that orbital, $I \approx -\varepsilon$.
\end{proposition}

\admitted

The approximation neglects the relaxation of the other electrons after the
ionisation; it is justified in the Year 3 volume. It turns a photoelectron
spectrum into an orbital diagram: water shows four valence bands, methane two,
in line with the \omterm{def:b2:fragment-orbitals:fragment}{fragment} diagrams --- and against the picture of two equivalent
lone pairs and two equivalent bonds in water, which describes the same total
electron density in a different way but does not give the energies of the
electrons removed.

\section{Ethene}

Ethene can be built from two \ce{CH2} \omterm{def:b2:fragment-orbitals:fragment}{fragments}. Each \ce{CH2} has, besides the
orbitals of its two C--H bonds, an orbital pointing away from the hydrogens in the
plane (which forms the C--C $\sigma$ bond with its partner) and a p orbital
perpendicular to the plane. The two p orbitals combine into a filled $\pi$ and an
empty $\pi^*$.

\begin{omfigure}
\begin{tikzpicture}[font=\small]
  \begin{scope}
    \pic at (0,0) {omp={0.85}{90}}; \pic at (1.2,0) {omp={0.85}{90}};
    \draw[thick] (0,0) -- (1.2,0);
    \draw[thick] (0,0) -- (-0.5,-0.45) (0,0) -- (-0.55,0.35) (1.2,0) -- (1.7,-0.45) (1.2,0) -- (1.75,0.35);
    \node at (0.6,-1.1) {$\pi$ (filled)};
  \end{scope}
  \begin{scope}[xshift=3.6cm]
    \pic at (0,0) {omp={0.85}{90}}; \pic at (1.2,0) {omp={-0.85}{90}};
    \draw[thick] (0,0) -- (1.2,0);
    \draw[thick] (0,0) -- (-0.5,-0.45) (0,0) -- (-0.55,0.35) (1.2,0) -- (1.7,-0.45) (1.2,0) -- (1.75,0.35);
    \node at (0.6,-1.1) {$\pi^*$ (empty)};
  \end{scope}
  \begin{scope}[xshift=7.4cm]
    \pic at (0,0) {omp={0.85}{90}}; \pic at (1.2,0) {omp={0.85}{0}};
    \draw[thick] (0,0) -- (1.2,0);
    \node at (0.6,-1.1) {twisted by $90^\circ$};
    \node[font=\scriptsize, align=center] at (0.6,1.3) {zero overlap};
  \end{scope}
\end{tikzpicture}

{\small The $\pi$ system of ethene, schematic (the molecule seen edge-on, the
hydrogens in the plane). The two p orbitals perpendicular to the plane combine
in phase ($\pi$, filled) and out of phase ($\pi^*$, empty). Twisting one \ce{CH2}
by $90^\circ$ makes the two p orbitals perpendicular: their overlap, and the $\pi$
bond, vanish.}
\end{omfigure}

The $\pi$ bond holds ethene flat: twisting one end turns the p orbitals away from
each other, their overlap falls as the cosine of the twist angle, and at $90^\circ$
the $\pi$ bond is gone. The energy of a $\pi$ bond must be supplied to twist the
molecule, which is why cis and trans isomers of alkenes do not interconvert at
room temperature, while rotation about a single bond is free.

\section{Hyperconjugation}

\begin{definition}[Hyperconjugation]\label{def:b2:fragment-orbitals:hyperconjugation}
\emph{Hyperconjugation}\index{hyperconjugation} is the interaction of a filled
$\sigma$ orbital of a C--H or C--C bond with a neighbouring empty or partly filled p
or $\pi^*$ orbital parallel to it.
\end{definition}

\begin{proposition}[Alkyl groups stabilise carbocations]\label{prop:b2:fragment-orbitals:hyperconjugation}
A carbocation is stabilised by each C--H or C--C bond on a neighbouring carbon that
can align with its empty p orbital: the more such bonds, the more stable the
cation, whence the order tertiary > secondary > primary > methyl.
\end{proposition}

\begin{proof}[Argument]
The empty p orbital of the cationic carbon and a filled $\sigma_{\text{C--H}}$ orbital
next to it form a two-level system with a large gap $\Delta$ and a small coupling
$\beta$ (non-zero only if the bond is roughly parallel to the p orbital). By
\cref{thm:b2:diatomic-mos:heteronuclear}, the filled $\sigma$ level is lowered by about
$\beta^2/\Delta$, and the two electrons gain twice that; the empty level rises, at no
cost. Each suitably aligned bond adds its own stabilisation; a methyl group
always has one C--H bond nearly aligned whatever its rotation.
\end{proof}

\begin{omfigure}
\begin{tikzpicture}[font=\small]
  \pic at (0,0) {omp={0.9}{90}};
  \node[circle, fill=black, inner sep=1.2pt] at (0,0) {};
  \node[font=\scriptsize, anchor=north west] at (0.08,-0.05) {C$^+$};
  \draw[thick] (0,0) -- (1.4,0);
  \node[circle, fill=black, inner sep=1.2pt] at (1.4,0) {};
  \draw[thick] (1.4,0) -- (1.75,0.85);
  \pic at (1.58,0.42) {omlobe={0.75}{68}};
  \pic at (1.58,0.42) {omlobe={0.45}{248}};
  \node[font=\scriptsize] at (1.95,1.0) {H};
  \draw[thick] (1.4,0) -- (2.0,-0.45) (1.4,0) -- (1.1,-0.6);
  \draw[<->, omProp, thick] (0.35,0.8) .. controls (0.8,1.2) and (1.2,1.2) .. (1.5,0.95);
  \node[font=\scriptsize, align=left, anchor=west] at (2.4,0.4) {filled $\sigma_{\text{C--H}}$ lobe parallel\\ to the empty p orbital};
  % level scheme
  \begin{scope}[xshift=7.2cm, yshift=-0.4cm]
    \pic (p) at (0,1.4) {omlevel={empty}{0.7}}; \node[left] at (-0.4,1.4) {empty p};
    \pic (s) at (2.2,-0.8) {omlevel={ud}{0.7}}; \node[right] at (2.6,-0.8) {$\sigma_{\text{C--H}}$};
    \pic (lo) at (1.1,-1.2) {omlevel={ud}{0.7}};
    \pic (hi) at (1.1,1.7) {omlevel={empty}{0.7}};
    \draw[omcorr, densely dashed] (p-r) -- (lo-l) (p-r) -- (hi-l) (s-l) -- (lo-r) (s-l) -- (hi-r);
    \node[font=\scriptsize, anchor=north] at (1.1,-1.35) {stabilised by $\approx \beta^2/\Delta$};
  \end{scope}
\end{tikzpicture}

{\small \omterm{def:b2:fragment-orbitals:hyperconjugation}{Hyperconjugation} in a carbocation, schematic. A filled C--H \omterm{def:b2:diatomic-mos:bonding}{bonding orbital}
on the neighbouring carbon, parallel to the empty p orbital, gives up part of its
electron density to it; the filled level is stabilised as in any interaction of
two levels of different energies.}
\end{omfigure}

\section{Exercises}

\begin{exercise}[$\star$]\label{exo:b2:fragment-orbitals:1}
Write the two \omterm{def:b2:fragment-orbitals:symmetry-adapted}{symmetry-adapted combinations} of the 1s orbitals of the hydrogens of
water and say which of them changes sign in the plane that bisects the H--O--H
angle perpendicular to the molecule.
\end{exercise}

\begin{exercise}[$\star$]\label{exo:b2:fragment-orbitals:2}
How many valence \omterm{def:b2:diatomic-mos:lcao}{molecular orbitals} does water have, how many valence electrons,
and how many filled valence orbitals?
\end{exercise}

\begin{exercise}[$\star$]\label{exo:b2:fragment-orbitals:3}
Why does the photoelectron spectrum of methane show two valence bands, one of them
three times as intense as the other?
\end{exercise}

\begin{exercise}[$\star$]\label{exo:b2:fragment-orbitals:4}
Why are the six atoms of ethene in one plane?
\end{exercise}

\begin{exercise}[$\star\star$]\label{exo:b2:fragment-orbitals:5}
The first ionisation energy of ammonia is \qty{10.07}{eV}, that of methane
\qty{12.61}{eV}. Which orbital loses the electron in each case, and why is ammonia the
easier to ionise?
\end{exercise}

\begin{exercise}[$\star\star$]\label{exo:b2:fragment-orbitals:6}
Compare the occupied levels of linear and bent water, and explain why a molecule
like \ce{BeH2}, with four valence electrons, is linear.
\end{exercise}

\begin{exercise}[$\star\star$]\label{exo:b2:fragment-orbitals:7}
Explain, with the two-level result, the order of stability of the carbocations
\ce{CH3+}, \ce{CH3CH2+}, \ce{(CH3)2CH+}, \ce{(CH3)3C+}.
\end{exercise}

\begin{exercise}[$\star\star$]\label{exo:b2:fragment-orbitals:8}
The overlap of two p orbitals twisted by an angle $\theta$ about the bond is
proportional to $\cos\theta$. With the \omterm{def:b2:diatomic-mos:integrals}{resonance integral} proportional to the
overlap, how does the $\pi$ stabilisation depend on $\theta$? At which angle does it
fall to half?
\end{exercise}

\begin{exercise}[$\star\star$]\label{exo:b2:fragment-orbitals:9}
Borane \ce{BH3} is planar, ammonia pyramidal. Using the \omterm{def:b2:fragment-orbitals:fragment}{fragment orbitals} of a
central atom and three hydrogens, explain the difference by the electrons each has.
\end{exercise}

\begin{exercise}[$\star\star\star$]\label{exo:b2:fragment-orbitals:10}
The \ce{H3+} ion is an equilateral triangle. With $\alpha$ and $\beta$ for the 1s
orbitals and overlap neglected, find its three orbital energies (the symmetric
combination has energy $\alpha + 2\beta$) and explain why two electrons suffice to
bind it.
\end{exercise}

\begin{exercise}[$\star\star\star$]\label{exo:b2:fragment-orbitals:11}
The allyl system \ce{CH2=CH-CH2} has three parallel p orbitals. Guess the shapes of
its three $\pi$ orbitals (signs on each carbon) from the number of nodes, and which
one is the highest occupied level of the allyl anion.
\end{exercise}

\begin{exercise}[$\star\star\star$]\label{exo:b2:fragment-orbitals:12}
Build the $\pi$ system of carbon dioxide from the 2p orbitals perpendicular to the
O--C--O axis: how many $\pi$ orbitals, which are bonding, nonbonding,
antibonding, and how many electrons do they hold?
\end{exercise}

\section{Problem: The Shape of Water}

\begin{problem}[{Weekend problem --- the oxygen and \ce{H2} \omterm{def:b2:fragment-orbitals:fragment}{fragments}, the orbitals of
a linear H--O--H, the levels that move when it bends, and the photoelectron
spectrum with Koopmans' approximation}]
\label{pb:b2:fragment-orbitals:1}
Data: geometry of water, $r(\ce{O-H}) = \qty{95.8}{pm}$, angle
$104.48^\circ$; first ionisation energies: \ce{H2O} \qty{12.62}{eV}, \ce{O}
\qty{13.62}{eV}. Water lies in the $yz$ plane, $z$ along the bisector.
% ledger: geo:H2O.rOH, geo:H2O.aHOH, ie:H2O, ie:O

\medskip
\noindent\textbf{Part I --- \omterm{def:b2:fragment-orbitals:fragment}{Fragments}.}
\begin{enumerate}
  \item How many valence electrons does water have, and from which atoms?
  \item Write the two \omterm{def:b2:fragment-orbitals:symmetry-adapted}{symmetry-adapted combinations} of the hydrogen 1s orbitals.
  \item Which oxygen orbitals are symmetric with respect to both mirror planes?
  \item Which oxygen orbital is antisymmetric with respect to the $xz$ plane only?
  \item Which oxygen orbital has no partner among the hydrogen combinations?
  \item How many \omterm{def:b2:diatomic-mos:lcao}{molecular orbitals} result, and how many are filled?
\end{enumerate}

\medskip
\noindent\textbf{Part II --- A linear H--O--H.}
\begin{enumerate}[resume]
  \item In a linear molecule along $z$, which combination interacts with the
        oxygen $2\mathrm p_z$?
  \item Which interacts with the oxygen 2s?
  \item Which oxygen orbitals stay nonbonding, and with which degeneracy?
  \item Fill the levels with the eight valence electrons.
  \item Would linear water be \omterm{def:b2:diatomic-mos:magnetism}{paramagnetic}?
  \item Where would its highest occupied level lie, compared with an oxygen 2p
        orbital?
\end{enumerate}

\medskip
\noindent\textbf{Part III --- Bending.}
\begin{enumerate}[resume]
  \item The linear molecule lying along $z$ is bent in the $yz$ plane. Which
        \omterm{def:b2:diatomic-mos:bonding}{nonbonding orbital} starts to overlap with $h_1 + h_2$?
  \item What happens to the filled level it carries?
  \item Which bonding level rises slightly, and why?
  \item Conclude: why is water bent?
  \item Compare the measured angle with $90^\circ$ (pure p bonding) and $109.5^\circ$.
  \item Which \omterm{def:b2:diatomic-mos:bonding}{nonbonding orbital} is left untouched by the bending?
\end{enumerate}

\medskip
\noindent\textbf{Part IV --- The photoelectron spectrum.}
\begin{enumerate}[resume]
  \item How many valence bands does the spectrum show?
  \item By Koopmans' approximation, from which orbital does the first band come?
  \item Why is that band narrow, while bands of \omterm{def:b2:diatomic-mos:bonding}{bonding orbitals} are broad?
  \item Compare the first ionisation energy of water with that of the oxygen
        atom, and explain the difference.
  \item What becomes of the ``two equivalent lone pairs'' of the Lewis structure?
  \item State the number of valence ionisation bands of water and the energy of
        the lowest.
\end{enumerate}
\end{problem}
